% Generated draft: requires mathematical and source review.
\documentclass[11pt]{article}
\usepackage[margin=25mm]{geometry}
\usepackage{fontspec}
\setmainfont{Latin Modern Roman}
\usepackage{amsmath,amssymb,mathtools,unicode-math}
\setmathfont{Latin Modern Math}
\usepackage{enumitem,needspace,xurl,hyperref}
\setlength{\parindent}{0pt}
\setlength{\parskip}{5pt}
\setlength{\emergencystretch}{3em}
\newcommand{\N}{\mathbb{N}}
\newcommand{\Z}{\mathbb{Z}}
\newcommand{\Q}{\mathbb{Q}}
\newcommand{\R}{\mathbb{R}}
\newcommand{\C}{\mathbb{C}}
\newcommand{\F}{\mathbb{F}}
\newcommand{\A}{\mathbb{A}}
\newcommand{\PP}{\mathbb{P}}
\newcommand{\E}{\mathbb{E}}
\newcommand{\eps}{\varepsilon}
\newcommand{\dd}{\,\mathrm{d}}
\begin{document}
\section*{UnsolvedMath ID: 701576}
\textbf{Source problem code:} catalogue entry
\subsection{Definitions and mathematical statement}
Let \(G\) be a group and let \(\varphi\in\operatorname{Aut}(G)\). Define \(\varphi\)-conjugacy by
\[
x\sim_\varphi y\quad\Longleftrightarrow\quad \exists z\in G\;\bigl(x=z^{-1}y\varphi(z)\bigr).
\]
This is an equivalence relation. Its number of equivalence classes, possibly infinite, is the Reidemeister number
\[
R(\varphi)=\lvert G/{\sim_\varphi}\rvert.
\]
The group \(G\) is finitely generated if \(G=\langle S\rangle\) for some finite subset \(S\subseteq G\). It is residually finite if, for every \(g\in G\setminus\{1\}\), there exist a finite group \(Q\) and a homomorphism \(q:G\to Q\) with \(q(g)\ne1\).

For a group \(H\), define \(H^{(0)}=H\) and \(H^{(n+1)}=[H^{(n)},H^{(n)}]\), where \([K,K]\) is generated by the commutators \(a^{-1}b^{-1}ab\) with \(a,b\in K\). The group \(H\) is soluble (solvable) if \(H^{(m)}=\{1\}\) for some integer \(m\ge0\). A group is virtually soluble, also called almost solvable in the supplied source, if it has a soluble subgroup of finite index.

\textbf{Conjecture (Fel'shtyn--Troitsky; Kourovka Notebook 19.28).} For every finitely generated residually finite group \(G\) and every \(\varphi\in\operatorname{Aut}(G)\),
\[
R(\varphi)<\infty
\quad\Longrightarrow\quad
\exists H\le G\;\exists m\ge0\;\bigl([G:H]<\infty\ \text{and}\ H^{(m)}=\{1\}\bigr).
\]

\textbf{Supplied background qualifications.} The source reports that the conjecture is proved when \(\varphi\) has prime order, citing Jabara (2008). It also reports that, if the finite generation hypothesis is dropped while retaining residual finiteness and the existence of an automorphism with finite Reidemeister number, virtual solubility need not follow, citing Dekimpe and Gonçalves (2014). It further reports a positive result for groups of finite upper rank, citing Troitsky (2025). In the usual terminology, finite upper rank means that there is an integer \(r\ge0\) such that every subgroup of every finite quotient of \(G\) can be generated by at most \(r\) elements. The supplied record does not give the full theorem statement for this last result, so its precise scope and any additional hypotheses require checking the cited paper. These background results and the supplied open status are source reports, not independent literature verification.
\subsection{Short English statement}
The conjecture asks whether every finitely generated residually finite group admitting an automorphism with finitely many twisted conjugacy classes has a soluble subgroup of finite index. Proposed by A. L. Fel'shtyn and E. V. Troitsky in 2018, it is labelled open in the supplied registry record; the accompanying special cases and qualifications are supplied literature reports.
\subsection{Sources}
\begin{itemize}
\item Local UnsolvedMath registry: permanent ID \(701576\), exact source problem code \texttt{catalogue entry}, title \emph{Kourovka Notebook Problem 19.28}. The supplied record attributes the problem to A. L. Fel'shtyn and E. V. Troitsky, with proposed year 2018. Its status is recorded as open; the record was updated on \texttt{2026-10-16T14:00:00Z}.
\item \emph{The Kourovka Notebook (Unsolved Problems in Group Theory)}, problem 19.28, 19th issue (2018); the supplied background also identifies the 21st issue edition (2026), arXiv:1401.0300v48. The registry's collection description names E. I. Khukhro and V. D. Mazurov as editors.
\item E. Jabara, \emph{J. Algebra} \textbf{320}, no. 10 (2008), 3671--3679. Cited by the supplied record for the prime-order automorphism case; no article title or link is supplied.
\item K. Dekimpe and D. Gonçalves, \emph{Bull. London Math. Soc.} \textbf{46}, no. 4 (2014), 737--746. Cited by the supplied record for failure of the conclusion without finite generation; no article title or link is supplied.
\item E. Troitsky, \emph{J. Group Theory} \textbf{28}, no. 1 (2025), 151--164. Cited by the supplied record for the finite-upper-rank result; no article title or link is supplied.
\item All attributions and literature assessments above are reproduced from the supplied record and have not been independently verified. No literature-verification date is supplied.
\end{itemize}
\end{document}
